\begin{pastexam}[label={px:nsysu-113-osds-03}]{中山資工所甲組 113 學年度 第 3 題}
（問答）

(10\%) Please indicate whether each of the following sorting algorithms is a stable sort or not:
bubble sort, heap sort, quick sort, merge sort, radix sort.
\end{pastexam}

\begin{pastsoln}
答案：bubble sort 穩定；heap sort 不穩定；quick sort 不穩定；merge sort 穩定；radix sort 穩定。

穩定的意思是：鍵相同的元素，排序後維持原本的前後次序（資結書第~14.1~節）。 「不穩定」只要舉出一個反例；
「穩定」要說明演算法為什麼不會讓相同的鍵越過對方。 下面用 $2_a, 2_b$ 表示兩個鍵都是 $2$、原本 $2_a$ 在前的
元素，各演算法都照課本的標準寫法。
\begin{itemize}
\item \textbf{bubble sort：穩定。} 只交換相鄰兩格，而且只在 $A[j] > A[j+1]$ 時才交換，相等的兩個元素
永遠不會交換，也就不會越過對方（資結書第~14.2~節）。
\item \textbf{heap sort：不穩定。} 反例 $\langle 2_a, 2_b \rangle$：建最大堆積時 $2_b$ 不比 $2_a$ 大，不交換；
接著把根 $2_a$ 和最後一格交換，得 $\langle 2_b, 2_a \rangle$，次序反了。 原因是把根換到陣列尾端是一次長距離的
交換。
\item \textbf{quick sort：不穩定。} 反例 $\langle 2_a, 2_b, 1 \rangle$，用取最後一個當基準的 Lomuto 分割
（資結書虛擬碼~14.8）：基準 $1$，$2_a, 2_b$ 都比它大，最後把基準和第一格交換，得 $\langle 1, 2_b, 2_a \rangle$，
之後的子陣列 $\langle 2_b, 2_a \rangle$ 已經排好不再變動。 分割時的遠距離交換會打亂相同鍵的次序。
\item \textbf{merge sort：穩定。} 合併時兩邊的鍵相同，先拿左半邊的（資結書虛擬碼~14.4 用 $\le$ 比較），左半
邊的元素原本就在右半邊之前，所以相同的鍵保持原次序。
\item \textbf{radix sort（LSD）：穩定。} 每一位數用穩定的計數排序，從最低位排到最高位；兩個鍵完全相同時，
每一趟都不會改變它們的次序（資結書第~15.2.2~節）。 LSD 基數排序的正確性本來就依賴每一趟都穩定。
\end{itemize}
\end{pastsoln}
